\documentclass[11pt]{article}

\usepackage[margin=0.8in]{geometry}
\usepackage{amsmath,amssymb,mathtools,array,longtable,enumitem,booktabs}
\usepackage{hyperref}

\setlist[itemize]{leftmargin=*}
\setlist[enumerate]{leftmargin=*}

\newcommand{\R}{\mathbb{R}}
\newcommand{\dd}{\,d}
\newcommand{\grad}{\nabla}
\newcommand{\curl}{\operatorname{curl}}
\newcommand{\divg}{\operatorname{div}}

\title{Quiz 5 Preparation Sheet\\Multivariable Calculus}
\author{Angad Singh Ahuja}
\date{Quiz: Friday 24/04/2026, 16:35--17:05}

\begin{document}
\maketitle

\tableofcontents
\newpage

\section*{How To Use This Sheet}
\addcontentsline{toc}{section}{How To Use This Sheet}

This sheet is built for a short quiz, so the priority is speed, recognition, and correct setup. The syllabus is: triple integrals in cylindrical coordinates, spherical coordinates, change of variables, vector fields, line integrals, the Fundamental Theorem for Line Integrals, Green's theorem, curl, and divergence.

Use the definitions first, then test yourself with the conceptual questions, then do the extracted problem-set questions. The short answer key is intentionally compact: it is meant for checking, not for replacing the full computations in your solution PDFs.

\section{Definitions And Core Facts}

\subsection{Triple Integrals}

A triple integral
\[
\iiint_E f(x,y,z)\,dV
\]
adds the values of a function over a three-dimensional solid $E$. If $f=1$, then
\[
\iiint_E 1\,dV
\]
is the volume of $E$.

For rectangular coordinates,
\[
dV=dx\,dy\,dz
\]
up to order, depending on the chosen order of integration.

\subsection{Cylindrical Coordinates}

Cylindrical coordinates are
\[
x=r\cos\theta,\qquad y=r\sin\theta,\qquad z=z.
\]
The radial coordinate is
\[
r=\sqrt{x^2+y^2}.
\]
The volume element is
\[
dV=r\,dz\,dr\,d\theta.
\]
The extra factor $r$ is the Jacobian.

Use cylindrical coordinates when the problem has cylinders, circular disks, cones of the form $z=\sqrt{x^2+y^2}$, paraboloids $z=a-br^2$, or expressions involving $x^2+y^2$.

\subsection{Spherical Coordinates}

Spherical coordinates are
\[
x=\rho\sin\phi\cos\theta,\qquad
y=\rho\sin\phi\sin\theta,\qquad
z=\rho\cos\phi.
\]
Here $\rho$ is distance from the origin, $\theta$ is the angle in the $xy$-plane, and $\phi$ is the angle from the positive $z$-axis.

Important identities:
\[
x^2+y^2+z^2=\rho^2,\qquad x^2+y^2=\rho^2\sin^2\phi,\qquad z=\rho\cos\phi.
\]
The volume element is
\[
dV=\rho^2\sin\phi\,d\rho\,d\phi\,d\theta.
\]

Use spherical coordinates for spheres, balls, spherical shells, cones through the origin, and integrands involving $x^2+y^2+z^2$.

\subsection{Change Of Variables In Double Integrals}

If
\[
x=x(u,v),\qquad y=y(u,v),
\]
then
\[
\iint_R f(x,y)\,dA
=\iint_S f(x(u,v),y(u,v))
\left|\frac{\partial(x,y)}{\partial(u,v)}\right|\,du\,dv.
\]
The Jacobian is
\[
\frac{\partial(x,y)}{\partial(u,v)}
=
\begin{vmatrix}
x_u & x_v\\
y_u & y_v
\end{vmatrix}.
\]

The main steps are:
\begin{enumerate}
\item Rewrite the boundary curves in terms of $u,v$.
\item Find the new region $S$.
\item Compute the Jacobian.
\item Replace the integrand and area element.
\item Integrate over the simpler region.
\end{enumerate}

\subsection{Vector Fields}

A vector field assigns a vector to each point. In two variables,
\[
\mathbf F(x,y)=\langle P(x,y),Q(x,y)\rangle.
\]
In three variables,
\[
\mathbf F(x,y,z)=\langle P(x,y,z),Q(x,y,z),R(x,y,z)\rangle.
\]

Vector fields can represent force, velocity, flow, electric fields, gradients, or direction fields.

\subsection{Line Integrals Of Scalar Functions}

For a scalar function $f$ along a curve $C$ parametrized by $\mathbf r(t)$, $a\le t\le b$,
\[
\int_C f\,ds=\int_a^b f(\mathbf r(t))\|\mathbf r'(t)\|\,dt.
\]

This adds the scalar value along the curve, weighted by arc length.

\subsection{Line Integrals Of Vector Fields}

For a vector field $\mathbf F$ along a parametrized curve $\mathbf r(t)$,
\[
\int_C \mathbf F\cdot d\mathbf r
=\int_a^b \mathbf F(\mathbf r(t))\cdot \mathbf r'(t)\,dt.
\]
In two dimensions, if $\mathbf F=\langle P,Q\rangle$, then
\[
\int_C \mathbf F\cdot d\mathbf r
=\int_C P\,dx+Q\,dy.
\]

This measures work done by the field along the curve.

\subsection{Conservative Vector Fields}

A vector field $\mathbf F$ is conservative if there exists a scalar potential function $\phi$ such that
\[
\mathbf F=\grad\phi.
\]
Then
\[
\int_C \mathbf F\cdot d\mathbf r
\]
depends only on the endpoints, not the path.

For $\mathbf F=\langle P,Q\rangle$ on a simply connected domain, a quick test is
\[
P_y=Q_x.
\]
If true, the field is conservative. If false, it is not conservative.

For $\mathbf F=\langle P,Q,R\rangle$ on a simply connected domain, a quick test is
\[
\curl \mathbf F=\mathbf 0.
\]

\subsection{Fundamental Theorem For Line Integrals}

If $\mathbf F=\grad\phi$, and $C$ goes from point $A$ to point $B$, then
\[
\int_C \mathbf F\cdot d\mathbf r=\phi(B)-\phi(A).
\]

This is the vector-field analogue of
\[
\int_a^b f'(x)\,dx=f(b)-f(a).
\]

\subsection{Green's Theorem}

Let $C$ be a positively oriented, simple, closed curve enclosing region $D$. If
\[
\mathbf F=\langle P,Q\rangle,
\]
then
\[
\oint_C P\,dx+Q\,dy
=\iint_D (Q_x-P_y)\,dA.
\]

Positive orientation means counterclockwise orientation.

Green's theorem converts a line integral around a closed curve into a double integral over the region inside the curve.

\subsection{Curl}

For a 3D vector field
\[
\mathbf F=\langle P,Q,R\rangle,
\]
the curl is
\[
\curl \mathbf F
=\grad\times \mathbf F
=
\left\langle
R_y-Q_z,\,
P_z-R_x,\,
Q_x-P_y
\right\rangle.
\]

Curl measures local rotation. If $\curl \mathbf F=\mathbf 0$ on a simply connected domain, then $\mathbf F$ is conservative.

For a 2D field $\mathbf F=\langle P,Q\rangle$, the scalar curl is
\[
Q_x-P_y.
\]
This is exactly the integrand in Green's theorem.

\subsection{Divergence}

For
\[
\mathbf F=\langle P,Q,R\rangle,
\]
the divergence is
\[
\divg \mathbf F=P_x+Q_y+R_z.
\]

Divergence measures net outward flow from a point. Positive divergence means the field behaves like a source; negative divergence means it behaves like a sink.

\section{Theoretical Short-Question Quiz}

\subsection*{A. Coordinates And Setup}
\begin{enumerate}
\item What is the difference between $r$ and $\rho$? \textbf{Answer:} $r=\sqrt{x^2+y^2}$ is distance from the $z$-axis; $\rho=\sqrt{x^2+y^2+z^2}$ is distance from the origin.
\item Why does cylindrical coordinates have Jacobian $r$? \textbf{Answer:} Cylindrical Jacobian is $r$ because polar area elements scale by radius.
\item Why does spherical coordinates have Jacobian $\rho^2\sin\phi$? \textbf{Answer:} Spherical Jacobian is $\rho^2\sin\phi$ because shells have radial scale $\rho^2$ and polar-angle scale $\sin\phi$.
\item What surfaces are represented by $r=a$, $\theta=c$, and $z=c$? \textbf{Answer:} $r=a$ is a cylinder; $\theta=c$ is a vertical half-plane; $z=c$ is a horizontal plane.
\item What surfaces are represented by $\rho=a$, $\phi=c$, and $\theta=c$? \textbf{Answer:} $\rho=a$ is a sphere; $\phi=c$ is a cone; $\theta=c$ is a vertical half-plane.
\item Convert $x^2+y^2=4$ into cylindrical coordinates. \textbf{Answer:} $x^2+y^2=4$ becomes $r=2$.
\item Convert $x^2+y^2+z^2=9$ into spherical coordinates. \textbf{Answer:} $x^2+y^2+z^2=9$ becomes $\rho=3$.
\item Convert $z=\sqrt{x^2+y^2}$ into spherical coordinates. \textbf{Answer:} $z=\sqrt{x^2+y^2}$ becomes $\phi=\pi/4$.
\item In cylindrical coordinates, why does $z=1-x^2-y^2$ become $z=1-r^2$? \textbf{Answer:} Since $x^2+y^2=r^2$, $z=1-x^2-y^2$ becomes $z=1-r^2$.
\item If a region is inside a sphere and above a cone, which coordinates are usually best? \textbf{Answer:} Spherical coordinates.
\item If a region is inside a cylinder and below a plane, which coordinates are usually best? \textbf{Answer:} Cylindrical coordinates.
\item What does the inequality $0\le \theta\le \pi/2$ usually mean geometrically? \textbf{Answer:} First quadrant in the $xy$-plane.
\item What does the inequality $0\le \phi\le \pi/4$ mean geometrically? \textbf{Answer:} Region near the positive $z$-axis above the cone $\phi=\pi/4$.
\item How do you decide whether $r$ or $z$ should be integrated first? \textbf{Answer:} Choose the order that makes bounds simplest; usually integrate the variable with simple top-bottom bounds first.
\item What is the biggest setup mistake students make in cylindrical/spherical coordinates? \textbf{Answer:} Forgetting the Jacobian.
\end{enumerate}

\subsection*{B. Change Of Variables}
\begin{enumerate}
\item What is a Jacobian? \textbf{Answer:} A Jacobian is the determinant measuring area/volume scaling under a coordinate change.
\item Why do we take the absolute value of the Jacobian in change of variables? \textbf{Answer:} Absolute value is used because area/volume is positive.
\item If the Jacobian is constant, what does that say geometrically? \textbf{Answer:} Constant Jacobian means constant area-scaling.
\item What is the first step when using a given change of variables? \textbf{Answer:} Transform the boundary.
\item Why should boundary curves be transformed before the integrand? \textbf{Answer:} Boundaries determine the new integration limits.
\item How do you know the new $uv$-region is simpler than the old $xy$-region? \textbf{Answer:} It usually turns curved boundaries into straight lines or circles.
\item What happens if the Jacobian is zero? \textbf{Answer:} The coordinate map degenerates and is not locally invertible.
\item Why is polar coordinates just a special case of change of variables? \textbf{Answer:} Polar coordinates are $x=r\cos\theta$, $y=r\sin\theta$.
\item In $x=r\cos\theta$, $y=r\sin\theta$, what is the Jacobian? \textbf{Answer:} The polar Jacobian is $r$.
\item If $x=2u$, $y=3v$, what is the Jacobian and what does it mean? \textbf{Answer:} For $x=2u$, $y=3v$, the Jacobian is $6$.
\end{enumerate}

\subsection*{C. Vector Fields And Line Integrals}
\begin{enumerate}
\item What is a vector field? \textbf{Answer:} A vector field assigns a vector to each point.
\item What does $\int_C \mathbf F\cdot d\mathbf r$ measure? \textbf{Answer:} $\int_C\mathbf F\cdot d\mathbf r$ measures work or circulation along a path.
\item Why does orientation matter for vector line integrals? \textbf{Answer:} Reversing orientation changes the sign.
\item Why does orientation not matter for scalar line integrals $\int_C f\,ds$? \textbf{Answer:} Scalar arc-length integrals use $ds$, which is always positive.
\item How do you parametrize a circle of radius $a$ counterclockwise? \textbf{Answer:} $\mathbf r(t)=(a\cos t,a\sin t)$, $0\le t\le2\pi$.
\item What is $d\mathbf r$ in terms of a parametrization $\mathbf r(t)$? \textbf{Answer:} $d\mathbf r=\mathbf r'(t)\,dt$.
\item When is a line integral path independent? \textbf{Answer:} A line integral is path independent when the field is conservative.
\item What is a conservative vector field? \textbf{Answer:} Conservative means $\mathbf F=\nabla\phi$.
\item How do you test whether $\langle P,Q\rangle$ is conservative in $\R^2$? \textbf{Answer:} In $\R^2$, check $P_y=Q_x$ on a simply connected domain.
\item How do you find a potential function once you know the field is conservative? \textbf{Answer:} Integrate one component, then compare partial derivatives to find the missing function.
\item Why is the domain important when checking if a field is conservative? \textbf{Answer:} Domains with holes can break global conservativeness.
\item What is special about $\left\langle \frac{-y}{x^2+y^2},\frac{x}{x^2+y^2}\right\rangle$? \textbf{Answer:} It has zero curl away from the origin but nonzero circulation around the origin.
\item Why can a field have zero curl but still fail to be globally conservative? \textbf{Answer:} Because the domain may not be simply connected (e.g., it has a hole).
\item What does the Fundamental Theorem for Line Integrals say? \textbf{Answer:} If $\mathbf F=\nabla\phi$, then $\int_C\mathbf F\cdot d\mathbf r=\phi(B)-\phi(A)$.
\item How is it similar to the one-variable Fundamental Theorem of Calculus? \textbf{Answer:} It is the multivariable version of $\int_a^b f'(x)\,dx=f(b)-f(a)$.
\end{enumerate}

\subsection*{D. Green's Theorem, Curl, Divergence}
\begin{enumerate}
\item State Green's theorem. \textbf{Answer:} Green's theorem: $\oint_C Pdx+Qdy=\iint_D(Q_x-P_y)dA$.
\item What does ``positively oriented'' mean? \textbf{Answer:} Positive orientation means counterclockwise.
\item In Green's theorem, why does $Q_x-P_y$ appear? \textbf{Answer:} $Q_x-P_y$ is scalar curl.
\item What is the scalar curl of $\langle P,Q\rangle$? \textbf{Answer:} Scalar curl is $Q_x-P_y$.
\item How can Green's theorem simplify a hard boundary integral? \textbf{Answer:} It replaces a boundary integral by an area integral.
\item What is the curl of a 3D vector field? \textbf{Answer:} $\curl\langle P,Q,R\rangle=\langle R_y-Q_z,P_z-R_x,Q_x-P_y\rangle$.
\item What does curl measure conceptually? \textbf{Answer:} Curl measures local rotation.
\item What is divergence? \textbf{Answer:} Divergence is $P_x+Q_y+R_z$.
\item What does divergence measure conceptually? \textbf{Answer:} Divergence measures source/sink behavior.
\item If $\mathbf F=\nabla\phi$, what is $\curl \mathbf F$? \textbf{Answer:} If $\mathbf F=\nabla\phi$, then $\curl\mathbf F=\mathbf0$.
\item If $\divg \mathbf F>0$, what does that suggest physically? \textbf{Answer:} Positive divergence suggests outward flow/source.
\item How do you remember the formula for curl? \textbf{Answer:} Remember curl as $\nabla\times\mathbf F$.
\item What is the difference between curl and divergence? \textbf{Answer:} Curl measures rotation; divergence measures expansion/outflow.
\item What is the difference between Green's theorem and the Fundamental Theorem for Line Integrals? \textbf{Answer:} FTLI is for conservative fields and endpoints; Green's theorem is for closed curves and enclosed regions.
\item Why must Green's theorem use a closed curve? \textbf{Answer:} Green's theorem needs a boundary of a region.
\end{enumerate}

\section{Direct Practice Questions From Problem Sets}

These are the most relevant questions to prepare for Quiz 5. They are copied or lightly standardized from Problem Sets 8, 9, and 10.

\subsection*{From Problem Set 8}
\begin{enumerate}
\item PS8 Problem 7(i): Evaluate $\iiint_V y\,dV$, where $V=\{0\le x\le 3,\ 0\le y\le x,\ x-y\le z\le x+y\}$.
\item PS8 Problem 7(iii): Evaluate $\iiint_V \sin y\,dV$, where $V$ is below $z=x$ and above the triangular region with vertices $(0,0,0)$, $(\pi,0,0)$, $(0,\pi,0)$.
\item PS8 Problem 7(iv): Evaluate $\iiint_V xy\,dV$, where $V$ is bounded by $y=x^2$, $x=y^2$, $z=0$, and $z=x+y$.
\item PS8 Problem 8(i): Change the order of $\int_0^1\int_0^y f(x,y)\,dx\,dy$.
\item PS8 Problem 8(ii): Change the order of $\int_0^2\int_{2y}^{y^2}f(x,y)\,dx\,dy$.
\item PS8 Problem 9(i): Write three other iterated integrals for the region $0\le x\le y\le z\le 1$.
\end{enumerate}

\subsection*{From Problem Set 9}
\begin{enumerate}
\item PS9 Problem 1(i): Use $x=u^2-v^2$, $y=2uv$ to evaluate $\iint_R y\,dA$, where $R$ is bounded by $y=0$, $y^2=4-4x$, and $y^2=4+4x$ with $y\ge0$.
\item PS9 Problem 1(ii): Use $x=2u+v$, $y=u+2v$ to evaluate $\iint_R(x-3y)\,dA$, where $R$ has vertices $(0,0)$, $(2,1)$, $(1,2)$.
\item PS9 Problem 2(ii): Evaluate $\int_{-2}^2\int_{-\sqrt{4-x^2}}^{\sqrt{4-x^2}}\int_{\sqrt{x^2+y^2}}^2(x^2+y^2)\,dz\,dy\,dx$.
\item PS9 Problem 4(i): Evaluate $\int_0^1\int_0^{\sqrt{1-x^2}}\int_{\sqrt{x^2+y^2}}^{\sqrt{2-x^2-y^2}}xy\,dz\,dy\,dx$ using spherical coordinates.
\item PS9 Problem 6(i): Find the volume above the cone $z=\sqrt{x^2+y^2}$ and below the sphere $x^2+y^2+z^2=z$.
\item PS9 Problem 7(ii): Determine whether $F=(e^x\sin y,e^x\cos y)$ is conservative. If yes, find $\phi$.
\item PS9 Problem 8(v): Determine whether $F=(e^x\sin yz,ze^x\cos yz,ye^x\cos yz)$ is conservative. If yes, find $\phi$.
\end{enumerate}

\subsection*{From Problem Set 10}
\begin{enumerate}
\item PS10 Problem 1(i): Show path independence and evaluate $\int_C 2xe^{-y}\,dx+(2y-x^2e^{-y})\,dy$ from $(1,0)$ to $(2,1)$.
\item PS10 Problem 1(ii): Show path independence and evaluate $\int_C\sin y\,dx+(x\cos y-\sin y)\,dy$ from $(2,0)$ to $(1,\pi)$.
\item PS10 Problem 2(i): Compute $\int_C \left\langle\frac{-y}{x^2+y^2},\frac{x}{x^2+y^2}\right\rangle\cdot d\mathbf r$ over the unit circle counterclockwise.
\item PS10 Problem 3(i): Use Green's theorem to evaluate $\int_C xy^2\,dx+2x^2y\,dy$ over the triangle with vertices $(0,0)$, $(2,2)$, $(2,4)$.
\item PS10 Problem 3(ii): Use Green's theorem to evaluate $\int_C\cos y\,dx+x^2\sin y\,dy$ over the rectangle with vertices $(0,0)$, $(5,0)$, $(5,2)$, $(0,2)$.
\item PS10 Problem 3(iii): Use Green's theorem to evaluate $\int_C x\,dy-y\,dx$ over the circle of radius $a$.
\end{enumerate}

\section{Short Answer Sheet}

\subsection*{Theory Quiz Short Answers}
\begin{enumerate}
\item $r=\sqrt{x^2+y^2}$ is distance from the $z$-axis; $\rho=\sqrt{x^2+y^2+z^2}$ is distance from the origin.
\item Cylindrical Jacobian is $r$ because polar area elements scale by radius.
\item Spherical Jacobian is $\rho^2\sin\phi$ because shells have radial scale $\rho^2$ and polar-angle scale $\sin\phi$.
\item $r=a$ is a cylinder; $\theta=c$ is a vertical half-plane; $z=c$ is a horizontal plane.
\item $\rho=a$ is a sphere; $\phi=c$ is a cone; $\theta=c$ is a vertical half-plane.
\item $x^2+y^2=4$ becomes $r=2$.
\item $x^2+y^2+z^2=9$ becomes $\rho=3$.
\item $z=\sqrt{x^2+y^2}$ becomes $\phi=\pi/4$.
\item Since $x^2+y^2=r^2$, $z=1-x^2-y^2$ becomes $z=1-r^2$.
\item Spherical coordinates.
\item Cylindrical coordinates.
\item First quadrant in the $xy$-plane.
\item Region near the positive $z$-axis above the cone $\phi=\pi/4$.
\item Choose the order that makes bounds simplest; usually integrate the variable with simple top-bottom bounds first.
\item Forgetting the Jacobian.
\item A Jacobian is the determinant measuring area/volume scaling under a coordinate change.
\item Absolute value is used because area/volume is positive.
\item Constant Jacobian means constant area-scaling.
\item Transform the boundary.
\item Boundaries determine the new integration limits.
\item It usually turns curved boundaries into straight lines or circles.
\item The coordinate map degenerates and is not locally invertible.
\item Polar coordinates are $x=r\cos\theta$, $y=r\sin\theta$.
\item The polar Jacobian is $r$.
\item For $x=2u$, $y=3v$, the Jacobian is $6$.
\item A vector field assigns a vector to each point.
\item $\int_C\mathbf F\cdot d\mathbf r$ measures work or circulation along a path.
\item Reversing orientation changes the sign.
\item Scalar arc-length integrals use $ds$, which is always positive.
\item $\mathbf r(t)=(a\cos t,a\sin t)$, $0\le t\le2\pi$.
\item $d\mathbf r=\mathbf r'(t)\,dt$.
\item A line integral is path independent when the field is conservative.
\item Conservative means $\mathbf F=\nabla\phi$.
\item In $\R^2$, check $P_y=Q_x$ on a simply connected domain.
\item Integrate one component, then compare partial derivatives to find the missing function.
\item Domains with holes can break global conservativeness.
\item It has zero curl away from the origin but nonzero circulation around the origin.
\item If $\mathbf F=\nabla\phi$, then $\int_C\mathbf F\cdot d\mathbf r=\phi(B)-\phi(A)$.
\item It is the multivariable version of $\int_a^b f'(x)\,dx=f(b)-f(a)$.
\item Green's theorem: $\oint_C Pdx+Qdy=\iint_D(Q_x-P_y)dA$.
\item Positive orientation means counterclockwise.
\item $Q_x-P_y$ is scalar curl.
\item Scalar curl is $Q_x-P_y$.
\item It replaces a boundary integral by an area integral.
\item $\curl\langle P,Q,R\rangle=\langle R_y-Q_z,P_z-R_x,Q_x-P_y\rangle$.
\item Curl measures local rotation.
\item Divergence is $P_x+Q_y+R_z$.
\item Divergence measures source/sink behavior.
\item If $\mathbf F=\nabla\phi$, then $\curl\mathbf F=\mathbf0$.
\item Positive divergence suggests outward flow/source.
\item Remember curl as $\nabla\times\mathbf F$.
\item Curl measures rotation; divergence measures expansion/outflow.
\item FTLI is for conservative fields and endpoints; Green's theorem is for closed curves and enclosed regions.
\item Green's theorem needs a boundary of a region.
\end{enumerate}

\subsection*{Selected Problem-Set Answers}

\begin{longtable}{>{\raggedright\arraybackslash}p{0.22\textwidth}>{\raggedright\arraybackslash}p{0.72\textwidth}}
\toprule
Question & Short answer \\
\midrule
PS8 7(i) & $\displaystyle \int_0^3\int_0^x\int_{x-y}^{x+y}y\,dz\,dy\,dx=\frac{27}{2}$.\\
PS8 7(iii) & $\displaystyle \frac12\int_0^\pi(\pi-y)^2\sin y\,dy=\frac{\pi^2-4}{2}$.\\
PS8 7(iv) & Region $0\le x\le1$, $x^2\le y\le\sqrt x$, $0\le z\le x+y$; value $\displaystyle \frac{3}{28}$.\\
PS8 8(i) & $\displaystyle \int_0^1\int_x^1 f(x,y)\,dy\,dx$.\\
PS8 8(ii) & Signed reversal: $\displaystyle \int_0^4\int_{\sqrt x}^{x/2}f(x,y)\,dy\,dx$. Geometric positive orientation: $-\int_0^4\int_{x/2}^{\sqrt x}f(x,y)\,dy\,dx$.\\
PS8 9(i) & Region $0\le x\le y\le z\le1$. Examples: $\int_0^1\int_0^z\int_x^z f\,dy\,dx\,dz$, $\int_0^1\int_x^1\int_x^z f\,dy\,dz\,dx$, $\int_0^1\int_x^1\int_y^1 f\,dz\,dy\,dx$.\\
PS9 1(i) & New region $0\le u,v\le1$, Jacobian $4(u^2+v^2)$, answer $2$.\\
PS9 1(ii) & New region $u\ge0$, $v\ge0$, $u+v\le1$, Jacobian $3$, answer $-3$.\\
PS9 2(ii) & Cylindrical setup $\int_0^{2\pi}\int_0^2\int_r^2 r^3\,dz\,dr\,d\theta$, answer $\frac{16\pi}{5}$.\\
PS9 4(i) & Spherical setup $0\le\theta\le\frac\pi2$, $0\le\phi\le\frac\pi4$, $0\le\rho\le\sqrt2$; answer $\frac{4\sqrt2-5}{15}$.\\
PS9 6(i) & $0\le\rho\le\cos\phi$, $0\le\phi\le\frac\pi4$, $0\le\theta\le2\pi$; volume $\frac{\pi}{8}$.\\
PS9 7(ii) & Conservative; potential $\phi=e^x\sin y+C$.\\
PS9 8(v) & Conservative; potential $\phi=e^x\sin(yz)+C$.\\
PS10 1(i) & Potential $\phi=x^2e^{-y}+y^2$; value $\frac4e$.\\
PS10 1(ii) & Potential $\phi=x\sin y+\cos y$; value $-2$.\\
PS10 2(i) & Parametrize unit circle; integral $2\pi$.\\
PS10 3(i) & Green's theorem gives $\iint_D2xy\,dA=\frac{56}{3}$.\\
PS10 3(ii) & Green's theorem gives $\int_0^5\int_0^2(2x+1)\sin y\,dy\,dx=30(1-\cos2)$.\\
PS10 3(iii) & Green's theorem gives $\iint_D2\,dA=2\pi a^2$.\\
\bottomrule
\end{longtable}

\section{Notation Table}

\begin{longtable}{>{\raggedright\arraybackslash}p{0.22\textwidth}>{\raggedright\arraybackslash}p{0.72\textwidth}}
\toprule
Notation & Meaning \\
\midrule
$\mathbf F$ & Vector field.\\
$P,Q,R$ & Component functions of a vector field.\\
$\grad\phi$ & Gradient of scalar potential $\phi$.\\
$\curl\mathbf F$ & Curl; measures local rotation.\\
$\divg\mathbf F$ & Divergence; measures source/sink behavior.\\
$C$ & Curve, often boundary curve.\\
$D$ & Plane region enclosed by $C$.\\
$E$ or $V$ & Solid region in $\R^3$.\\
$dA$ & Area element.\\
$dV$ & Volume element.\\
$ds$ & Arc-length element.\\
$d\mathbf r$ & Vector differential along a curve.\\
$r$ & Cylindrical/polar radius $\sqrt{x^2+y^2}$.\\
$\rho$ & Spherical radius $\sqrt{x^2+y^2+z^2}$.\\
$\theta$ & Angle in the $xy$-plane.\\
$\phi$ & Angle from positive $z$-axis in spherical coordinates.\\
$J$ & Jacobian determinant.\\
$\oint_C$ & Line integral around a closed curve.\\
$P_y=Q_x$ & Two-dimensional conservative-field test on simply connected domains.\\
$Q_x-P_y$ & Scalar curl; Green's theorem integrand.\\
\bottomrule
\end{longtable}

\section{FAQs}

\begin{enumerate}
\item \textbf{When should I use cylindrical coordinates?} Use them when $x^2+y^2$ appears, or the region has a cylinder, disk, cone, or circular paraboloid.
\item \textbf{When should I use spherical coordinates?} Use them for spheres, balls, cones, or integrands involving $x^2+y^2+z^2$.
\item \textbf{What is the most common mistake in cylindrical coordinates?} Forgetting the factor $r$ in $dV=r\,dz\,dr\,d\theta$.
\item \textbf{What is the most common mistake in spherical coordinates?} Forgetting $\rho^2\sin\phi$ or mixing up $\phi$ and $\theta$.
\item \textbf{Is $\phi$ measured from the $xy$-plane?} In Stewart convention, no. $\phi$ is measured from the positive $z$-axis.
\item \textbf{What does $z=\sqrt{x^2+y^2}$ mean geometrically?} A cone. In spherical coordinates it is $\phi=\pi/4$.
\item \textbf{How do I convert a sphere $x^2+y^2+z^2=z$?} Use $\rho^2=\rho\cos\phi$, so $\rho=\cos\phi$.
\item \textbf{Why does $x^2+y^2+z^2=z$ describe a sphere?} Completing the square gives $x^2+y^2+(z-\frac12)^2=\frac14$.
\item \textbf{Why is there an absolute value around the Jacobian?} Orientation can flip, but area/volume must stay positive.
\item \textbf{How do I find a new region after change of variables?} Transform each boundary curve into equations in the new variables.
\item \textbf{Do I transform the integrand or the region first?} Usually transform the region first.
\item \textbf{What does a conservative field mean?} It is a gradient field $\mathbf F=\nabla\phi$.
\item \textbf{Why are conservative fields important?} Their line integrals depend only on endpoints.
\item \textbf{How do I find a potential in 2D?} Integrate $P$ with respect to $x$, add $g(y)$, then compare with $Q$.
\item \textbf{How do I check if $\langle P,Q\rangle$ is conservative?} Check whether $P_y=Q_x$, assuming the domain is simply connected.
\item \textbf{What if $P_y=Q_x$ but the domain has a hole?} The field may fail to be globally conservative.
\item \textbf{What is the classic hole-domain example?} $\left\langle \frac{-y}{x^2+y^2},\frac{x}{x^2+y^2}\right\rangle$ on $\R^2\setminus\{0\}$.
\item \textbf{What is the line integral of that field around a counterclockwise circle?} $2\pi$, regardless of radius.
\item \textbf{Does zero curl always imply conservative?} Only on simply connected domains.
\item \textbf{What does Green's theorem do?} Converts a closed line integral into a double integral over the enclosed region.
\item \textbf{When should I use Green's theorem?} When the curve is closed and the boundary integral looks harder than the area integral.
\item \textbf{What is positive orientation?} Counterclockwise around the boundary.
\item \textbf{What happens if the curve is clockwise?} The answer changes sign.
\item \textbf{How do I remember Green's theorem?} $\oint_C Pdx+Qdy=\iint_D(Q_x-P_y)dA$.
\item \textbf{What if $Q_x-P_y$ is constant?} Then the line integral is constant times the area.
\item \textbf{Why is $\int_C x\,dy-y\,dx=2\cdot\text{Area}(D)$?} Here $P=-y$, $Q=x$, so $Q_x-P_y=1-(-1)=2$.
\item \textbf{What is curl in 2D?} The scalar $Q_x-P_y$.
\item \textbf{What is curl in 3D?} $\langle R_y-Q_z,P_z-R_x,Q_x-P_y\rangle$.
\item \textbf{What is divergence?} Sum of matching partial derivatives: $P_x+Q_y+R_z$.
\item \textbf{What is the difference between curl and divergence?} Curl measures rotation; divergence measures net outward flow.
\item \textbf{Do I need divergence theorem for this quiz?} The syllabus names divergence, but not necessarily the divergence theorem. Know the definition and interpretation.
\item \textbf{What is the fastest way to evaluate a conservative line integral?} Find $\phi$, then compute $\phi(\text{end})-\phi(\text{start})$.
\item \textbf{What if a line integral says ``any path''?} It is hinting at path independence/conservative fields.
\item \textbf{How do I parametrize a circle radius $a$?} $\mathbf r(t)=(a\cos t,a\sin t)$, $0\le t\le2\pi$.
\item \textbf{How do I parametrize clockwise?} Use $\mathbf r(t)=(a\cos t,-a\sin t)$ or reverse the limits.
\item \textbf{What is the answer if the vector field is a gradient and the curve is closed?} Zero.
\item \textbf{Why?} The start and end points are the same, so $\phi(B)-\phi(A)=0$.
\item \textbf{How do I decide bounds for a triple integral?} Draw/project the solid, choose outer variables for the projection, then set inner bounds between lower and upper surfaces.
\item \textbf{What is a good quiz strategy?} First identify the theorem/coordinate system, then write the correct setup before simplifying.
\item \textbf{What should I memorize?} Cylindrical formulas, spherical formulas, Jacobians, Green's theorem, FTLI, curl, divergence, and conservative-field tests.
\item \textbf{What should I practice most?} Setting up bounds and recognizing when a line integral is path independent.
\end{enumerate}

\end{document}

